\section{HHC can require infinitely many quadratic inequalities}
\label{sec:infinite-aggregation}

Fix $r\geq2$ and let
\begin{equation}\label{eq:ball-system}
 \begin{split}
 f_1(u,v)&=\|u\|^2-1,\qquad f_2(u,v)=\|v\|^2-1,\\
 f_3(u,v)&=\tfrac12-u\transpose v,\qquad
 S_r=\{(u,v)\in\R^r\times\R^r:f_1,f_2,f_3<0\}.
 \end{split}
\end{equation}
Write $C_r=\conv S_r$ and $D_r=\cl C_r$, and let $T_r$ be the
corresponding system with all three inequalities nonstrict.
The set $S_r$ is nonempty, since
$u=v=\sqrt{3/4}\,e_1$ is strictly feasible, and bounded by the two
unit-ball constraints. In particular, $C_r$ is nonempty and proper.

\begin{proposition}\label{prop:ball-good-cone}
The homogenization of~\eqref{eq:ball-system} has HHC. Its good
aggregation multipliers are exactly $\mathcal K_r\setminus\{0\}$, where
\begin{equation}\label{eq:ball-good-cone}
 \mathcal K_r=\{\lambda\in\R^3_+:\lambda_3^2\leq4\lambda_1\lambda_2\}.
\end{equation}
\end{proposition}
\begin{proof}
Apply Corollary~\ref{cor:gram-family} with $k=2$ and
$X=\begin{pmatrix}u\transpose\\v\transpose\end{pmatrix}$.
The homogeneous map is the linear image
$(G,T)\mapsto(G_{11}-T,G_{22}-T,T/2-G_{12})$ of
$\Phi_{2,r}$, so it has HHC.

Up to a coordinate permutation, the aggregated homogeneous matrix is
\begin{equation}\label{eq:ball-block-matrix}
 Q_\lambda=
 \left[\begin{pmatrix}
 \lambda_1&-\lambda_3/2\\-\lambda_3/2&\lambda_2
 \end{pmatrix}\otimes I_r\right]
 \oplus[\lambda_3/2-\lambda_1-\lambda_2].
\end{equation}
If its two-by-two block has a negative eigenvalue, that eigenvalue
repeats $r\geq2$ times; the aggregation cannot be good. Otherwise
the block is positive semidefinite, which for $\lambda\geq0$ is
equivalent to~\eqref{eq:ball-good-cone}. For a nonzero such multiplier,
\[
 \lambda_3\leq2\sqrt{\lambda_1\lambda_2}
 \leq\lambda_1+\lambda_2,
 \qquad \lambda_3/2-\lambda_1-\lambda_2<0.
\]
Thus $Q_\lambda$ has exactly one negative eigenvalue. Moreover,
$f_\lambda$ is convex and strictly negative on $S_r$; convexity
makes it strictly negative on every finite convex combination of
points in $S_r$. Both requirements for goodness hold.
\end{proof}

\subsection{A continuum of indispensable strict aggregations}

\begin{theorem}\label{thm:indispensable-rays}
Every exact description
\[
 C_r=\bigcap_{\lambda\in\Lambda}\{(u,v):f_\lambda(u,v)<0\},
 \qquad \Lambda\subseteq\mathcal K_r\setminus\{0\},
\]
contains a positive multiple of $(\tau,1/\tau,2)$ for every
$\tau\in[1,2]$. Hence such a description requires uncountably many
distinct rays. No finite family of nonstrict good aggregations
describes $D_r$.
\end{theorem}
\begin{proof}
Set $\epsilon=1/10$. For $\tau\in[1,2]$, choose $u_\tau,v_\tau$
with Gram matrix
\begin{equation}\label{eq:ray-witness-gram}
 G_\tau=
 \begin{pmatrix}1-\epsilon/\tau&2/5\\2/5&1-\epsilon\tau\end{pmatrix}.
\end{equation}
Both diagonal entries are at least $4/5$, and
$\det G_\tau\geq(4/5)^2-(2/5)^2=12/25>0$, so these vectors exist
in $\R^2$ and hence in $\R^r$. At $x_\tau=(u_\tau,v_\tau)$,
\begin{equation}\label{eq:ray-witness-slack}
 f(x_\tau)=\epsilon(-1/\tau,-\tau,1),\qquad
 f_\lambda(x_\tau)
 =\epsilon(-\lambda_1/\tau-\lambda_2\tau+\lambda_3)\leq0.
\end{equation}
The inequality follows from
$\lambda_1/\tau+\lambda_2\tau\geq2\sqrt{\lambda_1\lambda_2}
\geq\lambda_3$. For nonzero $\lambda\in\mathcal K_r$, equality
holds precisely when $\lambda$ is a positive multiple of
$\lambda(\tau)=(\tau,1/\tau,2)$. Indeed, equality forces both
positive diagonal weights, $\lambda_1/\tau=\lambda_2\tau$, and
$\lambda_3=2\sqrt{\lambda_1\lambda_2}$.

The strict valid aggregation $\lambda(\tau)$ excludes $x_\tau$
from $C_r$. Every other good ray is strictly satisfied there. Any
exact strict description must therefore include this ray, proving
the first assertion.

Now fix a finite good family and choose a ray $\lambda(\tau)$
it omits. Decrease both off-diagonal entries of $G_\tau$ by a
sufficiently small $\eta>0$. Positive definiteness is preserved,
and every member of the finite family remains strictly satisfied,
by continuity of its value as a linear function of the Gram entries.
The omitted aggregate now has value $2\eta>0$. It is nonpositive
on $D_r$, so the new Gram matrix realizes a point outside $D_r$
that satisfies every member of the finite family. This proves
the nonstrict assertion.
\end{proof}

Proposition~\ref{prop:ball-good-cone} and
Theorem~\ref{thm:indispensable-rays} give an affirmative answer to
Conjecture~3.1 of \citet{BDSpreprint}, already with three inequalities
in four variables. The indispensability proof does not use the general
HHC aggregation theorem. The exact description by all good aggregations
follows from the self-contained hull proof below and also from that
theorem.

\subsection{Ordinary and closed hulls, and a finite SDP lift}

For the rest of this section, abbreviate
\[
 p=1-\|u\|^2,\qquad q=1-\|v\|^2,\qquad h=u\transpose v.
\]
\begin{theorem}\label{thm:ball-hulls}
For every $r\geq2$,
\begin{align}
 C_r&=\{(u,v):p>0,\ q>0,\ h+\sqrt{pq}>1/2\},
       \label{eq:ball-open-hull}\\
 D_r=\conv T_r
    &=\{(u,v):p\geq0,\ q\geq0,\ h+\sqrt{pq}\geq1/2\}.
       \label{eq:ball-closed-hull}
\end{align}
Define the affine matrix
\begin{equation}\label{eq:ball-lift-matrix}
 M(u,v,\sigma)=
 \begin{pmatrix}
  1&\sigma&u\transpose\\
  \sigma&1&v\transpose\\
  u&v&I_r
 \end{pmatrix}.
\end{equation}
Then the two hulls have the finite lifted descriptions
\begin{align}
 C_r&=\{(u,v):\exists\sigma>1/2,\ M(u,v,\sigma)\succ0\},
          \label{eq:ball-strict-lift}\\
 D_r&=\{(u,v):\exists\sigma\geq1/2,\ M(u,v,\sigma)\succeq0\}.
          \label{eq:ball-closed-lift}
\end{align}
Furthermore, the coordinate multipliers $(1,0,0),(0,1,0)$ and
any countable dense family of $\lambda(\tau)=(\tau,1/\tau,2)$,
$\tau>0$, describe $D_r$ by nonstrict aggregations.
\end{theorem}
\begin{proof}
First we prove the strict hull formula directly. Necessity follows from
Proposition~\ref{prop:ball-good-cone}: its coordinate multipliers give
$p,q>0$ throughout $C_r$. At any such point, the nonzero good multiplier
$(q,p,2\sqrt{pq})$ has value
$2\sqrt{pq}(1/2-h-\sqrt{pq})$, which must be strictly negative.

For sufficiency, suppose $p,q>0$ and $h+\sqrt{pq}>1/2$.
Put $a=\sqrt p$ and $b=\sqrt q$. Since $r\geq2$, there is a unit
vector $e$ perpendicular to $bu-av$, including when this vector is zero.
Consequently
\[
 \gamma=\frac{u\transpose e}{a}=\frac{v\transpose e}{b}.
\]
Choose
\[
 \max\{0,(1/2-h)/(ab)\}<k<1,\qquad
 t_\pm=-\gamma\pm\sqrt{\gamma^2+k}.
\]
Such $k$ exists by the assumed strict scalar inequality.
The roots satisfy $t_-<0<t_+$ and $2\gamma t_\pm+t_\pm^2=k$.
For $x(t)=(u+tae,v+tbe)$, direct expansion at either root gives
\[
 \|u+tae\|^2=\|u\|^2+pk<1,\quad
 \|v+tbe\|^2=\|v\|^2+qk<1,\quad
 (u+tae)\transpose(v+tbe)=h+abk>1/2.
\]
Both points therefore belong to $S_r$, and
\[
 (u,v)=\frac{t_+}{t_+-t_-}x(t_-)
       +\frac{-t_-}{t_+-t_-}x(t_+)
\]
is a convex combination with positive weights. This proves
\eqref{eq:ball-open-hull} in every stated dimension, using only two
feasible points.

The cone is generated by the two coordinate rays and
$\lambda(\tau)$ for $\tau>0$. To check this directly, suppose
$\lambda_3>0$ and put $t=\lambda_3/2$. The cone inequality allows
a choice $t/\lambda_2\leq\tau\leq\lambda_1/t$. Then
\[
 \lambda=t(\tau,1/\tau,2)
       +(\lambda_1-t\tau)(1,0,0)
       +(\lambda_2-t/\tau)(0,1,0),
\]
with nonnegative coefficients. If $\lambda_3=0$, only the
coordinate rays are needed. Thus the strict intersection is given by
$p>0$, $q>0$, and
\begin{equation}\label{eq:ball-ray-inequality}
 \tau p+q/\tau>1-2h\qquad(\tau>0).
\end{equation}
For $p,q>0$, the left side attains its minimum $2\sqrt{pq}$ at
$\tau=\sqrt{q/p}$. This identifies the intersection of all strict good aggregations with
the already proved hull~\eqref{eq:ball-open-hull}.

The Schur complement of $I_r$ in~\eqref{eq:ball-lift-matrix} is
\[
 \begin{pmatrix}p&\sigma-h\\\sigma-h&q\end{pmatrix}.
\]
Consequently the projection in~\eqref{eq:ball-closed-lift} is
exactly the set on the right of~\eqref{eq:ball-closed-hull}:
a feasible $\sigma\geq1/2$ exists if and only if
$h+\sqrt{pq}\geq1/2$. The same calculation with strict positive
definiteness and $\sigma>1/2$ gives~\eqref{eq:ball-strict-lift}.
This uses that the interval $(h-\sqrt{pq},h+\sqrt{pq})$
intersects $(1/2,\infty)$ exactly when its upper endpoint exceeds
$1/2$.

Let $L_r$ denote the closed set on the right
of~\eqref{eq:ball-closed-hull}. It contains $C_r$, so $D_r\subseteq L_r$.
For $z=(u,v)\in L_r$, choose a feasible lift $\sigma$ and mix it
with the strictly positive definite matrix $M(0,0,1/2)$:
\[
 (1-\theta)M(u,v,\sigma)+\theta M(0,0,1/2)
 =M((1-\theta)u,(1-\theta)v,(1-\theta)\sigma+\theta/2)
 \succ0
\]
for $0<\theta<1$. Its scalar entry is at least $1/2$; if it
equals $1/2$, positive definiteness permits a sufficiently small
increase. Formula~\eqref{eq:ball-strict-lift} then puts
$(1-\theta)z$ in $C_r$. Sending $\theta\downarrow0$ proves
$L_r\subseteq D_r$.

The original closed set $T_r$ is compact, contains $S_r$, and
is contained in $L_r=D_r$. Its convex hull is compact and hence
closed, so
$D_r\subseteq\conv T_r\subseteq D_r$. This proves the asserted
equality for $T_r$ without assuming that closure commutes with
intersecting the original constraints.

Finally, the coordinate inequalities give $p,q\geq0$.
If the nonstrict version of~\eqref{eq:ball-ray-inequality} holds
on a dense set of positive $\tau$, continuity extends it to all
$\tau>0$. Its infimum is $2\sqrt{pq}$, also when $p=0$ or $q=0$
by taking a limit at an endpoint. The resulting condition is
exactly~\eqref{eq:ball-closed-hull}.
\end{proof}

The two-point construction is independent of the general HHC hull theorem
and works already at $r=2$. The choice of a direction perpendicular to
$bu-av$, with possibly unequal convex weights, avoids requiring a direction
perpendicular to both vectors. The latter midpoint construction would
require $r\geq3$.

\subsection{Cardinality limits for quadratic descriptions in the original variables}

The preceding obstruction concerned good aggregations. Even arbitrary
quadratics require uncountably many strict inequalities for the open
hull and infinitely many nonstrict inequalities for the closed hull.

\begin{theorem}\label{thm:no-finite-quadratics}
For $r\geq2$, no countable conjunction of strict quadratic inequalities
in $(u,v)$ describes $C_r$. No finite conjunction of nonstrict
quadratic inequalities in $(u,v)$ describes $D_r$. In both
statements the quadratics may be nonconvex and need not be
aggregations of the original three inequalities.
\end{theorem}
\begin{proof}
Restrict to the plane $u=xe_1$, $v=ye_2$. By
Theorem~\ref{thm:ball-hulls}, the planar sections are
\[
 \{(x,y):x^2<1,\ y^2<1,\ P(x,y)>0\},\qquad
 \{(x,y):x^2\leq1,\ y^2\leq1,\ P(x,y)\geq0\},
\]
respectively, where
\begin{equation}\label{eq:quartic-boundary}
 P(x,y)=(1-x^2)(1-y^2)-1/4.
\end{equation}
Both have the boundary arc
\[
 \Gamma=\{(x,\psi(x)):1/4\leq x\leq1/2\},\qquad
 \psi(x)=\sqrt{\frac{3/4-x^2}{1-x^2}}.
\]
The radicand is positive on a neighborhood of this interval and
is strictly below one. Hence $\Gamma$ is an analytic arc inside
the two ball bounds, with feasible and infeasible points in every
planar neighborhood.

We first show that a nonzero polynomial $q(x,y)$ of total degree
at most two can vanish at only finitely many points of $\Gamma$.
The rational function
\[
 R(x)=\frac{3/4-x^2}{1-x^2}
\]
is not a square in $\R(x)$, since its zeros at
$\pm\sqrt3/2$ and poles at $\pm1$ are simple. Thus
$y^2-R(x)$ is irreducible in $\R(x)[y]$. Divide $q$ by this
polynomial in $y$; the remainder has the form $a(x)y+b(x)$
with $a,b\in\R(x)$. If $q(x,\psi(x))$ vanished identically
on a subinterval and $a\ne0$, then
$\psi=-b/a$ there, forcing $R=(b/a)^2$ as a rational identity,
a contradiction. If $a=0$, then $b=0$ on that subinterval
and hence identically. In the remaining case $a=b=0$, the
polynomial $y^2-R$ divides $q$ in $\R(x)[y]$.

Viewed as a polynomial in $y$ over $\R[x]$, $P$ is primitive:
its two nonzero coefficients $x^2-1$ and $3/4-x^2$ are
coprime. Gauss's lemma therefore makes $P$ irreducible in
$\R[x,y]$ and transfers that divisibility to $P\mid q$ in
$\R[x,y]$. This is impossible for nonzero $q$ of degree at
most two, because $P$ has total degree four. We have shown
that $q(x,\psi(x))$ cannot vanish identically on a subinterval.
It is real analytic on a neighborhood of the compact parameter
interval; infinitely many zeros there would have an accumulation
point and force such identical vanishing. Its zeros on
$\Gamma$ are therefore finite.

Suppose a countable family of strict quadratic inequalities described
$C_r$, and restrict their polynomials to the plane. None is
identically zero, since the planar origin belongs to $C_r$ and
all its inequalities must be strict there. At a point of
$\Gamma$, continuity from the planar feasible region makes
all restricted polynomials nonpositive. At least one must
vanish, since otherwise the strict description would include
this boundary point. Thus a countable union of finite zero sets
would cover the uncountable arc $\Gamma$, a contradiction.

For a hypothetical finite nonstrict description of $D_r$, discard the
restricted polynomials that vanish identically on the plane.
Every remaining polynomial is nonpositive at each point of
$\Gamma$. At least one must vanish: if they were all strictly
negative, their finite conjunction would contain a planar
neighborhood, including points outside $D_r$. Finitely many finite
zero sets cannot cover $\Gamma$, giving a contradiction. The argument also covers the
case that no nonzero restrictions remain, since then the
description would include the entire plane.
\end{proof}

Theorem~\ref{thm:no-finite-quadratics} concerns conjunctions in
the original variables; the finite semidefinite lifts in
Theorem~\ref{thm:ball-hulls} remain available. For the closed hull,
the obstruction excludes only finite descriptions: a countable dense
family of nonstrict good aggregations suffices. These cardinality
statements give no runtime lower bound for optimization.

\subsection{Relation to previous infinite examples and finite bounds}

Infinite necessity itself predates the HHC framework.
\citet[Proposition~2.8 and Section~7.4]{DMS2022} give a closed
two-variable example with defining functions
\[
 g_1=x^2-1,\qquad g_2=y^2-1,\qquad
 g_3=-(x-1)^2-(y-1)^2+1.
\]
It fails HHC: on the homogeneous hyperplane $y=0$, the image is
$(x^2-t^2,-t^2,-x^2+2xt-t^2)$. The images of $(x,t)=(1,0)$
and $(0,1)$ have midpoint $(0,-1/2,-1)$, which would require
$x^2=t^2=1/2$ and $xt=0$. Moreover, their displayed
multipliers $(a^2,(1-a)^2,a^2-a+1)$, $0<a<1$, have leading
matrix $\operatorname{diag}(a-1,-a)$ and therefore at least
two negative homogeneous eigenvalues. That family is outside
the good-aggregation class used here. Another infinite
necessity result, \citet[Theorem~2]{DeyHanWang2026}, aggregates
boxed bilinear equalities with signed multipliers and convexifies
each aggregated set before intersecting. It concerns a
different aggregation operation.

In contrast, the construction in~\eqref{eq:ball-system} has HHC
and still requires a continuum of good rays for its strict hull, with
the stronger original-space obstruction in
Theorem~\ref{thm:no-finite-quadratics}. To the best of our knowledge,
this gives an explicit resolution of the HHC finiteness question in
\citet[Conjecture~3.1]{BDSpreprint}. We do not claim that infinite
aggregation in general, or a continuously varying family of active
inequalities, first arises here.

Semidefinite convexification is also established prior theory.
Quadratic matrix programming was studied by \citet{Beck2007}.
The convex-hull result of \citet[Section~4.1]{WangKilincKarzanPreprint}
applies to~\eqref{eq:ball-system} for $r\geq3$: take zero
objective, three constraints, replication count $r$, and the
positive definite sum $A_1+A_2=I_{2r}$ to meet their
definiteness assumption. Their paper appeared as
\citep{WangKilincKarzan2024}. This comparison concerns the
closed feasible-set hull, obtained from the zero-objective
epigraph, rather than an automatic replacement of nonstrict
inequalities by strict ones. \citet[Section~3.2.1, Corollary~1]{BrunSunWatson2026}
also derive a related inner-product hypograph hull over two
balls. Fixing the level after convexification does not by itself
prove the fixed-level hull here. The explicit proof above covers
$r=2$ and both choices of strictness; its finite lift is not
claimed as a new SDP convexification principle.

Finally, the example is outside the homogeneous definiteness
assumptions of known finiteness results. For any signed
$\lambda$, positive definiteness of the two-by-two block
in~\eqref{eq:ball-block-matrix} implies
$\lambda_1+\lambda_2>|\lambda_3|$, whereas positivity of the
scalar block requires
$\lambda_3/2>\lambda_1+\lambda_2$. These are incompatible.
Thus the $Q_i$ have no positive definite real linear combination,
even though $A_1+A_2\succ0$. The spectral determinant is
\[
 \det Q_\lambda=
 (\lambda_1\lambda_2-\lambda_3^2/4)^r
 (\lambda_3/2-\lambda_1-\lambda_2),
\]
which has a repeated factor. The common real projective zero set
is nonempty: take $t=1$, $u=e_1$, and
$v=e_1/2+\sqrt3\,e_2/2$. These facts prevent an application of
finiteness results requiring homogeneous PDLC, a smooth spectral
curve, or absence of common projective zeros. In particular, the
four-aggregation result of \citet[Theorem~1.4]{BDpreprint}, published
as \citep{BD2025}, has additional hypotheses and does not contradict
the construction here.
